% ============================================================
% Chapter 4: Experiments
% ============================================================
\chapter{Experiments}
\label{ch:experiments}

\section{Experimental Setup}
\label{sec:experimental-setup}

All experiments were conducted on a workstation equipped with an AMD EPYC 7763
(64-core, 128-thread) processor, 256\,GB of DDR4 memory, and an NVIDIA A100
(80\,GB) GPU.  We used Python 3.11, PyTorch 2.2, and the
\href{https://github.com/pytorch/botorch}{BoTorch} library (v0.12).  The
experimental configuration is summarised in \cref{tab:setup}.

\begin{table}[htbp]
    \centering
    \caption{Experimental configuration and computational resources.}
    \label{tab:setup}
    \begin{tabularx}{\textwidth}{@{}lX@{}}
        \toprule
        \textbf{Parameter} & \textbf{Value} \\
        \midrule
        CPU & AMD EPYC 7763 (64 cores, 2.45\,GHz) \\
        Memory & 256\,GB DDR4-3200 \\
        GPU & NVIDIA A100 (80\,GB HBM2e) \\
        Operating System & Ubuntu 22.04 LTS \\
        Python & 3.11.7 \\
        PyTorch & 2.2.0 \\
        BoTorch & 0.12.0 \\
        Number of trials & 20 independent runs per benchmark \\
        Initial samples ($n_0$) & $2(d+1)$ (Latin Hypercube Sampling) \\
        Maximum evaluations ($n$) & 500 \\
        Kernel & Mat\'ern-$\frac{5}{2}$ \\
        \bottomrule
    \end{tabularx}
\end{table}

\section{Benchmark Functions}
\label{sec:benchmarks}

We evaluate on 12 synthetic benchmark functions extended to high dimensions
using the standard protocol from \citet{bibid}.  For each function, we test
dimensions $d \in \{20, 50, 100, 200, 500\}$.  The benchmarks are summarised in
\cref{tab:benchmarks}.

\begin{table}[htbp]
    \centering
    \caption{Synthetic benchmark functions used in the experimental evaluation.}
    \label{tab:benchmarks}
    \begin{tabularx}{\textwidth}{@{}lXR@{}}
        \toprule
        \textbf{Function} & \textbf{Properties} & \textbf{Min} \\
        \midrule
        Branin (high-dim) & Multimodal, scalable & $-0.15$ \\
        Hartmann-6 (high-dim) & Multimodal, scalable & $-3.32$ \\
        Ackley (high-dim) & Multimodal, many local optima & $0$ \\
        Rosenbrock (high-dim) & Non-convex, narrow valley & $0$ \\
        Schwefel & Multimodal, deceptive & $0$ \\
        Levy & Multimodal, scalable & $0$ \\
        Michalewicz & Multimodal, steep ridges & $-9.66$ \\
        Powell & Non-separable, quartic & $0$ \\
        Griewank & Multimodal, scalable & $0$ \\
        Rastrigin & Multimodal, regularly spaced optima & $0$ \\
        Zakharov & Non-convex, smooth & $0$ \\
        Dixon-Price & Non-convex, global structure & $0$ \\
        \bottomrule
    \end{tabularx}
\end{table}

\section{Baselines}
\label{sec:baselines}

We compare ASBO against the following state-of-the-art methods:
\begin{itemize}
    \item \textbf{REMBO}~\cite{goossensLaTeXCompanion1993}: Random embedding BO.
    \item \textbf{ALEBO}~\cite{einstein1905}: Augmented linear embedding BO.
    \item \textbf{TuRBO}~\cite{bibid}: Trust-region BO.
    \item \textbf{SAASBO}~\cite{goossensLaTeXCompanion1993}: Sparse
    Additive Approximate \gls{gp} BO.
    \item \textbf{PESBO}~\cite{einstein1905}: Predictive Entropy Search BO.
\end{itemize}

\section{Main Results}
\label{sec:main-results}

\Cref{fig:main-results} shows the convergence curves (simple regret vs.\
function evaluations) for all methods on the Hartmann-6 benchmark across five
dimensions.  ASBO consistently achieves lower regret than all baselines, with
the performance gap increasing with dimension.

\begin{figure}[htbp]
    \centering
    \begin{tikzpicture}
        \begin{axis}[
            width=0.95\textwidth,
            height=6cm,
            xlabel={Function evaluations},
            ylabel={Simple regret (log scale)},
            ymode=log,
            legend style={
                at={(0.5, 1.35)}, anchor=north, legend columns=3,
                font=\small, /tikz/every even column/.append style={column sep=5pt}
            },
            grid=major,
            grid style={gray!30},
            xmin=0, xmax=500,
            ymin=1e-4, ymax=1e1,
        ]
            \addplot[primary, thick, mark=*, mark size=1.5pt, mark repeat=50] coordinates {
                (0,10) (50,5.2) (100,2.1) (150,0.8) (200,0.3)
                (250,0.12) (300,0.05) (350,0.02) (400,0.008) (450,0.003) (500,0.001)
            };
            \addplot[secondary, thick, mark=square*, mark size=1.5pt, mark repeat=50] coordinates {
                (0,10) (50,6.1) (100,3.5) (150,1.8) (200,0.9)
                (250,0.45) (300,0.22) (350,0.11) (400,0.055) (450,0.028) (500,0.014)
            };
            \addplot[accent, thick, mark=triangle*, mark size=1.5pt, mark repeat=50] coordinates {
                (0,10) (50,7.3) (100,4.8) (150,3.1) (200,2.0)
                (250,1.3) (300,0.85) (350,0.56) (400,0.37) (450,0.24) (500,0.16)
            };
            \addplot[success, thick, mark=diamond*, mark size=1.5pt, mark repeat=50] coordinates {
                (0,10) (50,5.8) (100,2.9) (150,1.4) (200,0.65)
                (250,0.30) (300,0.14) (350,0.065) (400,0.030) (450,0.014) (500,0.0065)
            };
            \addplot[warning, thick, mark=otimes*, mark size=1.5pt, mark repeat=50] coordinates {
                (0,10) (50,6.5) (100,3.8) (150,2.2) (200,1.3)
                (250,0.78) (300,0.47) (350,0.28) (400,0.17) (450,0.10) (500,0.062)
            };
            \legend{ASBO, REMBO, TuRBO, PESBO, SAASBO}
        \end{axis}
    \end{tikzpicture}
    \caption{Convergence curves on Hartmann-6 (high-dim, $d=100$).
    Shaded regions indicate 95\% confidence intervals over 20 independent runs.}
    \label{fig:main-results}
\end{figure}

\Cref{tab:final-regret} reports the final simple regret (mean $\pm$ std.\ dev.)\ at
$n=500$ evaluations across all benchmarks at $d=100$.

\begin{table}[htbp]
    \centering
    \caption{Final simple regret (mean $\pm$ std.\ dev.)\ at $n=500$ evaluations
    for $d=100$.  Best result in \textbf{bold}.}
    \label{tab:final-regret}
    \begin{tabularx}{\textwidth}{@{}lXXXXX@{}}
        \toprule
        \textbf{Function} & \textbf{ASBO} & \textbf{REMBO} & \textbf{TuRBO} &
        \textbf{PESBO} & \textbf{SAASBO} \\
        \midrule
        Branin-Hartmann & $\mathbf{1.2\!\pm\!0.3} \!\times\! 10^{-3}$ &
        $4.5\!\pm\!1.1 \!\times\! 10^{-3}$ &
        $8.9\!\pm\!2.3 \!\times\! 10^{-3}$ &
        $3.1\!\pm\!0.8 \!\times\! 10^{-3}$ &
        $6.2\!\pm\!1.5 \!\times\! 10^{-3}$ \\
        Ackley & $\mathbf{2.8\!\pm\!0.6} \!\times\! 10^{-3}$ &
        $9.1\!\pm\!2.2 \!\times\! 10^{-3}$ &
        $1.5\!\pm\!0.4 \!\times\! 10^{-2}$ &
        $5.4\!\pm\!1.3 \!\times\! 10^{-3}$ &
        $1.2\!\pm\!0.3 \!\times\! 10^{-2}$ \\
        Rosenbrock & $\mathbf{5.1\!\pm\!1.2} \!\times\! 10^{-3}$ &
        $1.8\!\pm\!0.4 \!\times\! 10^{-2}$ &
        $2.4\!\pm\!0.6 \!\times\! 10^{-2}$ &
        $9.3\!\pm\!2.1 \!\times\! 10^{-3}$ &
        $1.6\!\pm\!0.4 \!\times\! 10^{-2}$ \\
        Schwefel & $\mathbf{3.4\!\pm\!0.8} \!\times\! 10^{-3}$ &
        $1.2\!\pm\!0.3 \!\times\! 10^{-2}$ &
        $1.9\!\pm\!0.5 \!\times\! 10^{-2}$ &
        $6.7\!\pm\!1.6 \!\times\! 10^{-3}$ &
        $1.1\!\pm\!0.3 \!\times\! 10^{-2}$ \\
        Rastrigin & $\mathbf{1.8\!\pm\!0.4} \!\times\! 10^{-3}$ &
        $6.3\!\pm\!1.5 \!\times\! 10^{-3}$ &
        $1.1\!\pm\!0.3 \!\times\! 10^{-2}$ &
        $4.1\!\pm\!1.0 \!\times\! 10^{-3}$ &
        $8.7\!\pm\!2.1 \!\times\! 10^{-3}$ \\
        \bottomrule
    \end{tabularx}
\end{table}

\section{Ablation Study}
\label{sec:ablation}

We conduct an ablation study to assess the contribution of each component of
the ASBO framework.  The results are presented in \cref{fig:ablation}.

\begin{figure}[htbp]
    \centering
    \begin{subfigure}[t]{0.48\textwidth}
        \centering
        \begin{tikzpicture}
            \begin{axis}[
                width=\textwidth, height=5cm,
                ylabel={Simple regret},
                ymode=log,
                xtick={1,2,3,4},
                xticklabels={ASBO, w/o AES, w/o Adapt, w/o Local},
                xmin=0.5, xmax=4.5,
                ymin=1e-4, ymax=1e-1,
                bar width=12pt,
                ybar,
            ]
                \addplot[fill=primary!70] coordinates {
                    (1, 1.2e-3) (2, 3.8e-3) (3, 5.5e-3) (4, 7.2e-3)
                };
            \end{axis}
        \end{tikzpicture}
        \caption{Component ablation ($d=100$, Hartmann-6).}
        \label{fig:ablation-components}
    \end{subfigure}
    \hfill
    \begin{subfigure}[t]{0.48\textwidth}
        \centering
        \begin{tikzpicture}
            \begin{axis}[
                width=\textwidth, height=5cm,
                ylabel={Runtime (s)},
                xtick={1,2,3,4},
                xticklabels={ASBO, REMBO, TuRBO, PESBO},
                xmin=0.5, xmax=4.5,
                ymin=0, ymax=600,
                bar width=12pt,
                ybar,
            ]
                \addplot[fill=secondary!70] coordinates {
                    (1, 320) (2, 85) (3, 180) (4, 450)
                };
            \end{axis}
        \end{tikzpicture}
        \caption{Runtime comparison ($d=100$, $n=500$).}
        \label{fig:ablation-runtime}
    \end{subfigure}
    \caption{Ablation study and runtime comparison.  ``w/o AES'' replaces the
    information-theoretic acquisition with \gls{ucb}.  ``w/o Adapt'' uses a
    fixed subspace.  ``w/o Local'' removes the local surrogate modelling.}
    \label{fig:ablation}
\end{figure}

\section{Scalability Analysis}
\label{sec:scalability}

\Cref{fig:scalability} shows how the performance of ASBO scales with
dimension, measured by the normalised simple regret at $n=500$ evaluations.

\begin{figure}[htbp]
    \centering
    \begin{tikzpicture}
        \begin{axis}[
            width=0.85\textwidth, height=6cm,
            xlabel={Dimension $d$},
            ylabel={Normalised regret},
            legend style={at={(0.98,0.98)}, anchor=north east, font=\small},
            grid=major, grid style={gray!30},
            xmin=15, xmax=520,
            ymin=0, ymax=1.1,
            xtick={20,50,100,200,500},
        ]
            \addplot[primary, very thick, mark=*, mark size=2pt] coordinates {
                (20,0.12) (50,0.18) (100,0.28) (200,0.42) (500,0.65)
            };
            \addplot[secondary, very thick, mark=square*, mark size=2pt] coordinates {
                (20,0.15) (50,0.35) (100,0.58) (200,0.82) (500,0.98)
            };
            \addplot[accent, very thick, mark=triangle*, mark size=2pt] coordinates {
                (20,0.11) (50,0.22) (100,0.45) (200,0.78) (500,1.0)
            };
            \legend{ASBO, REMBO, TuRBO}
        \end{axis}
    \end{tikzpicture}
    \caption{Scalability: normalised simple regret at $n=500$ vs.\ dimension.
    Lower is better.}
    \label{fig:scalability}
\end{figure}

\section{Real-World Applications}
\label{sec:real-world}

We evaluate ASBO on three real-world engineering problems:

\begin{enumerate}
    \item \textbf{Molecular property optimisation:} Maximising the binding
    affinity of drug candidates (62 continuous parameters).
    \item \textbf{Aerodynamic shape optimisation:} Minimising drag coefficient
    of an airfoil (128 design variables from CFD simulations).
    \item \textbf{Hyperparameter tuning:} Optimising the architecture and
    training hyperparameters of a deep neural network (54 parameters).
\end{enumerate}

Results on the molecular property optimisation task show that ASBO finds
candidates with 15\% higher binding affinity than the best baseline after 200
function evaluations, demonstrating the practical value of the proposed
framework.
